\section{Results}
\label{sec:results}

\subsection{Measured telemetry adds value beyond cycle and operating state}

Table~\ref{tab:input-comparison} gives the nested outer-test results. Elapsed
cycle alone is a strong baseline, with AUROC 0.923 and final-window recall
0.563 at a 0.050 pre-window false-alert rate. Adding operating conditions
raises AUROC to 0.935 but leaves recall at 0.554. The all-observable model,
which adds summaries from the 14 measured channels, reaches AUROC 0.966 and
recall 0.793 at a 0.051 false-alert rate.

The paired outer-test comparison of all observable inputs against cycle plus
operating conditions estimates a 0.031 AUROC gain, with a 95\% engine-cluster
interval from 0.023 to 0.039. Its recall gain is 0.239, with an interval from
0.183 to 0.296. The false-alert-rate difference is 0.004, with an interval
that spans zero (-0.006 to 0.015). At the matched alert budget, the added
telemetry therefore increases the share of final-window cycles found without a
clear increase in the pre-window alert rate.

The telemetry-only rule is also useful without a cycle-number feature. It has
AUROC 0.939, precision 0.863, and recall 0.771 at a 0.047 false-alert rate.
Compared with cycle only, it gains 0.208 recall (95\% interval 0.105 to 0.310).
Its AUROC difference is 0.016, with an interval from -0.009 to 0.038. The
recall result is the stronger evidence for the fixed 5\% warning policy.

\begin{table}[H]
\centering
\caption{Nested engine-held-out performance at a target 5\% pre-window
false-alert budget. Values in brackets are 95\% percentile intervals from
2,000 complete-engine bootstrap resamples.}
\label{tab:input-comparison}
\scriptsize
\resizebox{\textwidth}{!}{%
\begin{tabular}{lrrrr}
\toprule
Input & AUROC & Precision & Final-window recall & Pre-window false-alert rate \\
\midrule
Elapsed cycle only & 0.923 [0.910, 0.939] & 0.813 [0.766, 0.866] & 0.563 [0.480, 0.646] & 0.050 [0.032, 0.067] \\
Cycle plus operating conditions & 0.935 [0.922, 0.947] & 0.821 [0.775, 0.869] & 0.554 [0.480, 0.629] & 0.046 [0.031, 0.063] \\
Measured telemetry only & 0.939 [0.918, 0.957] & 0.863 [0.837, 0.888] & 0.771 [0.715, 0.824] & 0.047 [0.037, 0.058] \\
All observable inputs & 0.966 [0.958, 0.973] & 0.857 [0.821, 0.892] & 0.793 [0.751, 0.835] & 0.051 [0.037, 0.066] \\
\bottomrule
\end{tabular}}
\end{table}

\begin{figure}[H]
\centering
\includegraphics[width=0.98\linewidth]{maintenance_telemetry_benchmark}
\caption{Measured telemetry improves final-window classification beyond
elapsed cycle and operating conditions. Each bar uses only outer-test engine
predictions; error bars are 95\% complete-engine bootstrap intervals.}
\label{fig:input-comparison}
\end{figure}

\subsection{Nested selection improves the all-observable warning policy}

The information comparison above keeps the classifier fixed. To test whether a
stronger warning policy is available without using test-engine labels, the
secondary all-observable analysis selects from four tabular candidates inside
each outer-training split. Histogram gradient boosting is selected in three
outer folds and the RBF support-vector classifier in two. The resulting policy
reaches AUROC 0.973, precision 0.861, and final-window recall 0.821 at a 0.051
realized false-alert rate (Table~\ref{tab:model-selection}). It warns 54 of 90
engines at or before the final-window start, so its gain is in cycle-level
detection rather than first-warning lead time.

Relative to the fixed all-observable random forest, nested model selection adds
0.007 AUROC and 0.028 final-window recall at the same pooled false-alert rate.
The paired complete-engine bootstrap interval is 0.002 to 0.011 for AUROC and
0.004 to 0.051 for recall. The model-family choice and threshold are made from
inner out-of-fold predictions in every outer fold, so these intervals compare
outer-test predictions only.

\begin{table}[H]
\centering
\caption{All-observable warning policies at a target 5\% pre-window
false-alert budget. The selected policy chooses its model family and threshold
within each outer-training split. Values in brackets are 95\% complete-engine
bootstrap intervals.}
\label{tab:model-selection}
\scriptsize
\resizebox{\textwidth}{!}{%
\begin{tabular}{lrrrr}
\toprule
Policy & AUROC & Precision & Final-window recall & Pre-window false-alert rate \\
\midrule
Fixed random forest & 0.966 [0.958, 0.973] & 0.857 [0.821, 0.892] & 0.793 [0.751, 0.835] & 0.051 [0.037, 0.066] \\
Nested selected policy & 0.973 [0.966, 0.979] & 0.861 [0.827, 0.894] & 0.821 [0.778, 0.862] & 0.051 [0.038, 0.066] \\
\bottomrule
\end{tabular}}
\end{table}

\begin{figure}[H]
\centering
\includegraphics[width=0.98\linewidth]{maintenance_nested_model_selection}
\caption{Nested model selection improves the all-observable policy without
using labels from test engines. The selected policy uses histogram gradient
boosting in three outer folds and an RBF support-vector classifier in two.}
\label{fig:model-selection}
\end{figure}

\subsection{The fixed 0.5 threshold is a useful diagnostic, not an equivalent policy}

Table~\ref{tab:fixed-threshold} evaluates the same outer-test scores at a
fixed 0.5 threshold. The telemetry-only rule changes little because its nested
thresholds range from 0.460 to 0.479. The all-observable, cycle-only, and
cycle-plus-operating policies require higher thresholds to meet the 5\% alert
budget. Their 0.5-threshold recalls are higher, but so are their false-alert
rates. This confirms that threshold selection materially affects the decision
policy and was not omitted from the primary result.

\begin{table}[H]
\centering
\caption{Diagnostic fixed-threshold performance on the same outer-test
probabilities. These values do not replace nested threshold selection.}
\label{tab:fixed-threshold}
\small
\begin{tabular}{lrrr}
\toprule
Input & Precision & Final-window recall & Pre-window false-alert rate \\
\midrule
Elapsed cycle only & 0.624 & 0.887 & 0.204 \\
Cycle plus operating conditions & 0.700 & 0.859 & 0.141 \\
Measured telemetry only & 0.890 & 0.751 & 0.036 \\
All observable inputs & 0.791 & 0.872 & 0.089 \\
\bottomrule
\end{tabular}
\end{table}

\subsection{Sensor-only scoring produces earlier first warnings}

The first-warning analysis evaluates a different operational property from
cycle-level recall: whether a policy flags an engine at or before the start of
the final 20-cycle window. The telemetry-only rule warns 69 of 90 engines by
that point, with a median lead of 6 cycles among the 87 engines it eventually
warns (Table~\ref{tab:warning-timing}). The fixed all-observable random forest
warns 50 engines by the window start, while the nested selected policy warns
54. Both have a median lead of 1 cycle among the engines they eventually warn.

This difference is a tradeoff, not a contradiction. The all-observable model
was selected for cycle-level final-window classification, while the first-alert
metric rewards an earlier threshold crossing. A deployment study that needs
lead time should select and validate an engine-level early-warning objective
directly.

\begin{table}[H]
\centering
\caption{First-warning timing from outer-test predictions. Positive lead means
the first warning occurs before the final-window start.}
\label{tab:warning-timing}
\small
\begin{tabular}{lrrrr}
\toprule
Input & By window start & Late warning & No warning & Median lead among warned \\
\midrule
Measured telemetry only & 69/90 (76.7\%) & 18/90 & 3/90 & 6 cycles \\
All observable, fixed random forest & 50/90 (55.6\%) & 39/90 & 1/90 & 1 cycle \\
All observable, nested selected policy & 54/90 (60.0\%) & 34/90 & 2/90 & 1 cycle \\
\bottomrule
\end{tabular}
\end{table}

\begin{figure}[H]
\centering
\includegraphics[width=0.98\linewidth]{maintenance_sensor_warning_timing}
\caption{First warning for the telemetry-only policy across 90 outer-test
engine trajectories. This policy uses no cycle number or operating-condition
field. Positive lead values arrive before the final 20-cycle window begins.}
\label{fig:warning-timing}
\end{figure}

\subsection{The conclusion is stable across model seeds}

Table~\ref{tab:seed-robustness} reports five complete nested evaluations with
random-forest seeds 2024 through 2028. The seed ranges are narrow for each
input. They do not replace external validation, because the simulated engines
and the outer folds remain the same.

\begin{table}[H]
\centering
\caption{Five-seed robustness check. Each entry is mean [minimum, maximum]
over full nested evaluations using fixed outer engine splits. The first four
rows use the fixed random-forest reference; the final row repeats model-family
selection inside every outer-training split.}
\label{tab:seed-robustness}
\small
\begin{tabular}{lrrr}
\toprule
Input & AUROC & Final-window recall & Pre-window false-alert rate \\
\midrule
Elapsed cycle only & 0.923 [0.923, 0.923] & 0.564 [0.559, 0.568] & 0.050 [0.048, 0.051] \\
Cycle plus operating conditions & 0.935 [0.934, 0.935] & 0.559 [0.554, 0.566] & 0.049 [0.046, 0.049] \\
Measured telemetry only & 0.940 [0.939, 0.941] & 0.774 [0.771, 0.776] & 0.048 [0.046, 0.050] \\
All observable inputs & 0.965 [0.965, 0.966] & 0.785 [0.780, 0.793] & 0.050 [0.049, 0.051] \\
Nested selected policy & 0.973 [0.973, 0.973] & 0.821 [0.821, 0.822] & 0.051 [0.051, 0.051] \\
\bottomrule
\end{tabular}
\end{table}

\FloatBarrier
